\documentclass[dvipdfmx,11pt,aspectratio=169]{beamer}

% 日本語対応
\usepackage{bxdpx-beamer}
\usepackage{pxjahyper}
\renewcommand{\kanjifamilydefault}{\gtdefault}

% テーマ設定
% Google Slidesのようなモダンでシンプルな見た目にするため metropolis を採用
\usetheme{metropolis}
\metroset{block=fill}
\setbeamertemplate{navigation symbols}{}

% 色設定 (Google Slidesテーマから抽出)
\definecolor{GoogleBlue}{HTML}{4285F4}
\definecolor{GoogleRed}{HTML}{EA4335}
\definecolor{GoogleYellow}{HTML}{FBBC05}
\definecolor{GoogleGreen}{HTML}{34A853}
\definecolor{GoogleDarkGray}{HTML}{212121}
\definecolor{GoogleLightGray}{HTML}{EEEEEE}

\setbeamercolor{normal text}{fg=GoogleDarkGray, bg=white}
\setbeamercolor{alerted text}{fg=GoogleRed}
\setbeamercolor{example text}{fg=GoogleGreen}

\setbeamercolor{palette primary}{fg=white, bg=GoogleBlue}
\setbeamercolor{palette secondary}{fg=white, bg=GoogleRed}
\setbeamercolor{palette tertiary}{fg=white, bg=GoogleYellow}
\setbeamercolor{palette quaternary}{fg=white, bg=GoogleGreen}

\setbeamercolor{title separator}{fg=GoogleBlue}
\setbeamercolor{progress bar}{fg=GoogleBlue, bg=GoogleLightGray}
\setbeamercolor{frametitle}{fg=white, bg=GoogleBlue}

% フォント設定
\usepackage[T1]{fontenc}
\usepackage{lmodern}

% 図解用
\usepackage{tikz}
\usetikzlibrary{shapes, arrows.meta, positioning}

% タイトル情報
\title{物理的制約のあるレガシーRANを次世代コアへ統合する\\疎結合アーキテクチャの提案と評価}
\subtitle{〜世代間プロトコル変換による移行コスト削減と機能拡張〜}
\author{Taihei}
\institute{所属大学/学部}
\date{\today}

\begin{document}

\begin{frame}
  \titlepage
\end{frame}

\begin{frame}{目次}
  \tableofcontents
\end{frame}

\section{はじめに}
\begin{frame}{背景：CNの進化とRAN更新のギャップ}
  \begin{itemize}
    \item \textbf{コアネットワークの進化 (ソフトウェア化)}
      \begin{itemize}
        \item CNは5G以降でソフトウェア化されたため高速に進化
        \item 6Gでも5Gコアをベースに継続的な機能拡張が予想される
      \end{itemize}
    \item \textbf{現状のデプロイメントとその限界}
      \begin{itemize}
        \item 5Gコア特有の新機能(スライシング等)を活用するには、SA (Standalone) 構成へ移行する必要がある
        \item 現在はRANを先行更新するNSA (Non-Standalone)構成が主流だが、CNは4Gのままであり新機能が使えない
        \item SA化には対応RANへの更新が必要だが、コストや物理的制約(例: NTN)により困難な場合がある
      \end{itemize}
    $\Rightarrow$ \textbf{Coreの進化速度に対し、RANの更新サイクルにギャップが生じる}
  \end{itemize}
\end{frame}

\begin{frame}{NSA構成とSA構成の比較}
  \centering
  \begin{tikzpicture}[scale=0.8, every node/.style={font=\small}]
    % NSA構成（左側）
    \node[font=\normalsize\bfseries] at (2.5,4.5) {NSA (Non-Standalone)};

    % UE
    \node[draw, rounded corners, fill=GoogleLightGray, minimum width=1.5cm, minimum height=0.8cm] (nsa-ue) at (0,2) {UE};

    % eNB
    \node[draw, rounded corners, fill=GoogleBlue!50, minimum width=1.5cm, minimum height=0.8cm] (nsa-enb) at (2.5,3) {eNB};
    % gNB
    \node[draw, rounded corners, fill=GoogleGreen!50, minimum width=1.5cm, minimum height=0.8cm] (nsa-gnb) at (2.5,1) {gNB};

    % EPC
    \node[draw, rounded corners, fill=GoogleBlue!50, minimum width=1.5cm, minimum height=0.8cm] (nsa-epc) at (5,2) {4G EPC};

    % 接続線
    \draw[->, thick] (nsa-ue) -- (nsa-enb);
    \draw[->, thick] (nsa-ue) -- (nsa-gnb);
    \draw[->, thick] (nsa-enb) -- (nsa-epc);
    \draw[->, thick, dashed] (nsa-gnb) -- (nsa-enb);

    % 説明
    \node[text width=4cm, align=center] at (2.5,-0.5) {5G RANを追加するが\\Coreは4Gのまま};

    % 区切り線
    \draw[thick, dashed, gray] (7,5) -- (7,-1);

    % SA構成（右側）
    \node[font=\normalsize\bfseries] at (11,4.5) {SA (Standalone)};

    % UE
    \node[draw, rounded corners, fill=GoogleLightGray, minimum width=1.5cm, minimum height=0.8cm] (sa-ue) at (8.5,2) {UE};

    % gNB
    \node[draw, rounded corners, fill=GoogleGreen!50, minimum width=1.5cm, minimum height=0.8cm] (sa-gnb) at (11,2) {gNB};

    % 5GC
    \node[draw, rounded corners, fill=GoogleGreen!50, minimum width=1.5cm, minimum height=0.8cm] (sa-5gc) at (13.5,2) {5GC};

    % 接続線
    \draw[->, thick] (sa-ue) -- (sa-gnb);
    \draw[->, thick] (sa-gnb) -- (sa-5gc);

    % 説明
    \node[text width=4cm, align=center] at (11,-0.5) {5G RANと5GCの組み合わせ\\新機能が利用可能};
  \end{tikzpicture}
\end{frame}

\begin{frame}{課題：ギャップが問題になるケース}
  \begin{itemize}
    \item \textbf{研究開発：新機能の先行検証}
      \begin{itemize}
        \item 次世代CN機能（スライシング、URLLC等）を検証したい
        \item 対応RAN（5G gNB）は高価で免許取得も必要
        \item $\Rightarrow$ \textbf{低コストかつ手軽に検証できる環境が必要}
      \end{itemize}
    \item \textbf{商用導入：マイグレーション戦略の柔軟性欠如}
      \begin{itemize}
        \item 商用網を前提にした場合、全RANを一斉に更新するのは現実的でない
        \item 段階移行のためのOptionもあるが、移行期間中は旧RANエリアで新CN機能が利用できない
        \item $\Rightarrow$ \textbf{RAN更新を伴わないCN先行移行パスが必要}
      \end{itemize}
    \item \textbf{運用維持：物理的に更新が難しいRAN}
      \begin{itemize}
        \item NTNでは衛星やHAPSを使用するため、一度配備すると地上局よりも物理的な更新が困難
        \item 地上のCNが進化しても、空のRANは旧世代のまま残る
        \item $\Rightarrow$ \textbf{新旧世代が共存可能な収容技術が必要}
      \end{itemize}
  \end{itemize}
\end{frame}

\begin{frame}{NTN (Non-Terrestrial Network) とは}
  \centering
  \begin{tikzpicture}[scale=0.75, every node/.style={font=\small}]
    % 地上エリア
    \fill[GoogleLightGray!50] (-0.5,-0.5) rectangle (12,0.8);
    \node at (0.5,0.15) {\textbf{地上}};

    % 地上局/ゲートウェイ
    \node[draw, rounded corners, fill=GoogleBlue!50, minimum width=1.8cm, minimum height=0.7cm] (gw) at (1,1.8) {地上局};

    % 地上のCore
    \node[draw, rounded corners, fill=GoogleGreen!50, minimum width=1.8cm, minimum height=0.7cm] (core) at (1,3.1) {5GC/6GC};
    \draw[->, thick] (gw) -- (core);

    % 衛星（宇宙）
    \node[draw, circle, fill=GoogleYellow!50, minimum size=1.2cm] (sat) at (5,5) {衛星};
    \node[above] at (5,5.7) {\textbf{宇宙}};
    \node[below, text width=2cm, align=center] at (5,4.2) {\scriptsize 4G RAN搭載};

    % HAPS（成層圏）
    \node[draw, ellipse, fill=GoogleYellow!50, minimum width=1.8cm, minimum height=0.7cm] (haps) at (9,3.5) {HAPS};
    \node[above] at (9,4.1) {\textbf{成層圏}};
    \node[below, text width=2cm, align=center] at (9,2.9) {\scriptsize 4G RAN搭載};

    % UE（地上）
    \node[draw, rounded corners, fill=white, minimum width=1.2cm, minimum height=0.5cm] (ue1) at (5,0.15) {UE};
    \node[draw, rounded corners, fill=white, minimum width=1.2cm, minimum height=0.5cm] (ue2) at (8,0.15) {UE};

    % 接続線
    \draw[->, thick, dashed] (sat) -- (gw) node[midway, above, sloped] {\scriptsize Feeder};
    \draw[->, thick] (sat) -- (ue1) node[midway, right] {\scriptsize Service};
    \draw[->, thick] (haps) -- (ue2);
    \draw[->, thick, dashed] (haps) -- (gw);

    % 問題点（右側に配置）
    \node[draw, rounded corners, fill=GoogleRed!20, text width=4cm, align=center, font=\small] at (10.5,1.5) {打ち上げ後はHW更新不可\\$\Rightarrow$ SW更新のみ対応可能};
  \end{tikzpicture}
\end{frame}

\begin{frame}{目的とアプローチ}
  \begin{block}{目的}
    物理的に更新困難なレガシーRANを、ハードウェア変更なしに次世代Coreに収容し、段階的な移行と新機能利用を可能にするアーキテクチャを提案する。
  \end{block}
  \begin{itemize}
    \item \textbf{アプローチ：プロトコル変換による疎結合化}
      \begin{itemize}
        \item 変換ゲートウェイを介在させ、RANとCoreの直接的な依存を排除
        \item 世代間で不変な機能（AAA, Mobility, U-plane）を抽象化してマッピング
      \end{itemize}
    \item \textbf{本研究のスコープ}
      \begin{itemize}
        \item 4G eNB (LTE) と 5G Core (5GC) の異世代相互接続を実証
      \end{itemize}
    \item \textbf{期待される効果（課題への回答）}
      \begin{itemize}
        \item \textbf{R\&D:} 既存RANを用いた安価な次世代Core検証環境
        \item \textbf{移行:} RAN更新を伴わない柔軟なCore先行移行
        \item \textbf{運用:} NTNなど物理制約のあるインフラの延命と機能拡張
      \end{itemize}
  \end{itemize}
\end{frame}

\section{提案手法}
\begin{frame}{提案アーキテクチャ：S1-N2/N3 変換ゲートウェイ}
  \centering
  \begin{tikzpicture}[scale=0.9, every node/.style={font=\small}]
    % Nodes
    \node[draw, rounded corners, fill=GoogleBlue!20, minimum width=2cm, minimum height=1cm] (enb) at (0,0) {4G eNB};
    \node[draw, rounded corners, fill=GoogleGreen!20, minimum width=2cm, minimum height=1cm] (core) at (10,0) {5G Core};

    % Converter Box
    \node[draw, thick, fill=GoogleYellow!10, minimum width=4.5cm, minimum height=2.5cm, dashed] (conv) at (5,0) {};
    \node[above] at (5,1.3) {\textbf{提案コンバータ}};

    % Internal Logic
    \node[draw, fill=white, minimum width=1.2cm] (s1) at (3.8,0) {S1};
    \node[draw, fill=white, minimum width=1.2cm] (ng) at (6.2,0) {NG};
    \draw[<->, thick] (s1) -- (ng) node[midway, above] {\scriptsize 変換};

    % Connections
    \draw[<->, thick] (enb) -- (s1) node[midway, above] {S1};
    \draw[<->, thick] (ng) -- (core) node[midway, above] {N2/N3};

    % Roles
    \node[below=0.2cm of enb, align=center, font=\scriptsize] {EPCだと思って\\接続};
    \node[below=0.2cm of core, align=center, font=\scriptsize] {gNBだと思って\\接続};
    \node[below=1.4cm of conv, align=center, font=\scriptsize] {相互にエミュレーションを行い\\ギャップを吸収};
  \end{tikzpicture}
\end{frame}

\begin{frame}{世代不変性に基づく疎結合化の原理}
  \begin{block}{キーアイデア}
    モバイル通信の根幹機能(AAA, Mobility, QoS)は世代を超えて\alert{本質的に同じ}であり、プロトコル表現が異なるだけ $\Rightarrow$ 抽象化して相互変換可能
  \end{block}
  \begin{itemize}
    \item \textbf{C-Plane (制御信号): ステートマシン変換}
      \begin{itemize}
        \item 接続管理(Attach/Registration)、Mobility(HO)等の状態遷移は世代不変
        \item S1-AP $\leftrightarrow$ NG-AP のメッセージ名・パラメータをマッピング
      \end{itemize}
    \item \textbf{U-Plane (データ転送): QoS抽象化}
      \begin{itemize}
        \item Bearer (4G) も QoS Flow (5G) も「優先度付きトンネル」という点で同一
        \item GTP-U TEID とQoS識別子 (EBI $\leftrightarrow$ QFI) を付け替え
      \end{itemize}
    \item \textbf{NAS (端末-Core通信): 透過転送}
      \begin{itemize}
        \item \alert{NAS over LTE:} 5G NASを4GのNASコンテナに格納して運ぶ
        \item これにより端末は4G無線上で5G Coreへ直接登録可能(技術的なブレイクスルー)
      \end{itemize}
  \end{itemize}
\end{frame}

\section{設計}
\begin{frame}{システム構成：プロトコルスタック}
  \centering
  \begin{tikzpicture}[scale=0.75, every node/.style={font=\small}]
    % eNB側
    \node[font=\normalsize\bfseries] at (1.5,5.5) {eNB};
    \node[draw, fill=GoogleBlue!20, minimum width=2.5cm, minimum height=0.6cm] at (1.5,4.5) {S1-AP};
    \node[draw, fill=GoogleBlue!20, minimum width=2.5cm, minimum height=0.6cm] at (1.5,3.8) {SCTP};
    \node[draw, fill=GoogleBlue!20, minimum width=2.5cm, minimum height=0.6cm] at (1.5,3.1) {IP (172.24.0.111)};

    % GTP-U
    \node[draw, fill=GoogleBlue!20, minimum width=2.5cm, minimum height=0.6cm] at (1.5,2.0) {GTP-U};
    \node[draw, fill=GoogleBlue!20, minimum width=2.5cm, minimum height=0.6cm] at (1.5,1.3) {UDP:2152};

    % Converter左側
    \node[font=\normalsize\bfseries] at (5.5,5.5) {S1N2 Converter};
    \node[draw, fill=GoogleYellow!20, minimum width=2.5cm, minimum height=0.6cm] at (5.5,4.5) {S1-AP};
    \node[draw, fill=GoogleYellow!20, minimum width=2.5cm, minimum height=0.6cm] at (5.5,3.8) {SCTP:36412};

    % 変換ロジック
    \node[draw, thick, fill=GoogleRed!10, minimum width=2.5cm, minimum height=1.2cm, align=center] at (5.5,2.6) {変換層\\{\scriptsize S1↔NG}\\{\scriptsize TEID管理}};

    % GTP Proxy
    \node[draw, fill=GoogleYellow!20, minimum width=2.5cm, minimum height=0.6cm] at (5.5,1.3) {GTP Proxy};

    % Converter右側
    \node[draw, fill=GoogleYellow!20, minimum width=2.5cm, minimum height=0.6cm] at (9.5,4.5) {NG-AP};
    \node[draw, fill=GoogleYellow!20, minimum width=2.5cm, minimum height=0.6cm] at (9.5,3.8) {SCTP:38412};
    \node[draw, fill=GoogleYellow!20, minimum width=2.5cm, minimum height=0.6cm] at (9.5,3.1) {IP (172.24.0.30)};

    % AMF/UPF側
    \node[font=\normalsize\bfseries] at (13,5.5) {AMF};
    \node[draw, fill=GoogleGreen!20, minimum width=2.5cm, minimum height=0.6cm] at (13,4.5) {NG-AP};
    \node[draw, fill=GoogleGreen!20, minimum width=2.5cm, minimum height=0.6cm] at (13,3.8) {SCTP};
    \node[draw, fill=GoogleGreen!20, minimum width=2.5cm, minimum height=0.6cm] at (13,3.1) {IP (172.24.0.12)};

    \node[font=\normalsize\bfseries] at (13,1.8) {UPF};
    \node[draw, fill=GoogleGreen!20, minimum width=2.5cm, minimum height=0.6cm] at (13,1.0) {GTP-U};
    \node[draw, fill=GoogleGreen!20, minimum width=2.5cm, minimum height=0.6cm] at (13,0.3) {IP (172.24.0.21)};

    % 接続線 C-Plane
    \draw[<->, thick] (2.75,4.5) -- (4.25,4.5);
    \draw[<->, thick] (6.75,4.5) -- (8.25,4.5);
    \draw[<->, thick] (10.75,4.5) -- (11.75,4.5);

    % 接続線 U-Plane
    \draw[<->, thick, dashed] (2.75,1.3) -- (4.25,1.3);
    \draw[<->, thick, dashed] (6.75,1.3) -- (11.75,1.0);

    % ラベル
    \node[font=\scriptsize] at (3.5,4.8) {C-Plane};
    \node[font=\scriptsize] at (3.5,1.6) {U-Plane};
  \end{tikzpicture}
\end{frame}

\begin{frame}{UE接続シーケンス：S1 Setup}
  \centering
  \begin{tikzpicture}[scale=0.7, every node/.style={font=\small}]
    % Actors
    \node[font=\normalsize\bfseries] at (0,6) {eNB};
    \node[font=\normalsize\bfseries] at (5,6) {S1N2 Converter};
    \node[font=\normalsize\bfseries] at (10,6) {AMF};

    % Lifelines
    \draw[thick] (0,5.5) -- (0,0);
    \draw[thick] (5,5.5) -- (5,0);
    \draw[thick] (10,5.5) -- (10,0);

    % Messages
    \draw[->, thick, GoogleBlue] (0,5) -- (5,4.8);
    \node[above, font=\scriptsize, align=left] at (2.5,5) {S1SetupRequest\\{\tiny Global-ENB-ID, TAC}};

    \draw[->, thick, GoogleGreen] (5,4.5) -- (10,4.3);
    \node[above, font=\scriptsize, align=left] at (7.5,4.5) {NGSetupRequest\\{\tiny Global-gNB-ID, TAC}};

    \draw[<-, thick, GoogleGreen] (5,3.8) -- (10,3.6);
    \node[below, font=\scriptsize, align=left] at (7.5,3.6) {NGSetupResponse\\{\tiny AMF Name, GUAMI}};

    \draw[<-, thick, GoogleBlue] (0,3.3) -- (5,3.1);
    \node[below, font=\scriptsize, align=left] at (2.5,3.1) {S1SetupResponse\\{\tiny MME Name}};

    % 変換処理
    \node[draw, fill=GoogleYellow!20, text width=3cm, align=center, font=\tiny] at (5,1.2) {ID変換\\ENB-ID → gNB-ID\\TAC保持};
  \end{tikzpicture}
\end{frame}

\begin{frame}{完全シーケンス(1/5)：Initial UE Message + NAS変換}
  \centering
  \begin{tikzpicture}[scale=0.55, every node/.style={font=\tiny}]
    % Actors
    \node[font=\small\bfseries] at (0,9) {UE};
    \node[font=\small\bfseries] at (4,9) {eNB};
    \node[font=\small\bfseries] at (8,9) {S1N2};
    \node[font=\small\bfseries] at (12,9) {AMF};

    % Lifelines
    \draw[thick] (0,8.5) -- (0,0);
    \draw[thick] (4,8.5) -- (4,0);
    \draw[thick] (8,8.5) -- (8,0);
    \draw[thick] (12,8.5) -- (12,0);

    % 1. RRC Connection
    \draw[->, thick] (0.1,8) -- (3.9,8);
    \node[above, font=\tiny] at (2,8.1) {RRC Connection Request};

    \draw[->, thick] (3.9,7.2) -- (0.1,7.2);
    \node[above, font=\tiny] at (2,7.3) {RRC Connection Setup};

    % 2. Attach Request (4G NAS)
    \draw[->, thick, GoogleBlue] (0.1,6.2) -- (3.9,6.2);
    \node[above, font=\tiny] at (2,6.3) {RRC Setup Cmp + Attach Req};

    % NAS詳細 (4G) - 左側に配置
    \node[draw, fill=GoogleBlue!10, text width=2.3cm, align=left, font=\tiny, anchor=north east] at (-0.3,5.9) {
      \textbf{4G NAS: Attach Req}\\
      - IMSI\\
      - UE Network Cap\\
      - PDN Conn Req
    };

    % 3. S1AP InitialUEMessage
    \draw[->, thick, GoogleBlue] (4.1,4.2) -- (7.9,4.2);
    \node[above, font=\tiny] at (6,4.3) {InitialUEMessage (S1AP)};

    % 4. NAS変換処理 - 中央下に配置
    \node[draw, thick, fill=GoogleYellow!30, text width=2.6cm, align=center, font=\tiny, anchor=north] at (8,3.6) {
      \textbf{NAS変換}\\
      Attach → Registration\\
      IMSI → SUCI
    };

    % 5. NGAP InitialUEMessage
    \draw[->, thick, GoogleGreen] (8.1,2.0) -- (11.9,2.0);
    \node[above, font=\tiny] at (10,2.1) {InitialUEMessage (NGAP)};

    % NAS詳細 (5G) - 右側に配置
    \node[draw, fill=GoogleGreen!10, text width=2.3cm, align=left, font=\tiny, anchor=north west] at (12.3,1.7) {
      \textbf{5G NAS: Reg Req}\\
      - SUCI\\
      - UE Security Cap\\
      - Requested NSSAI
    };
  \end{tikzpicture}
\end{frame}

%% ========================================
%% シーケンス 2/5: Authentication（詳細）
%% ========================================
\begin{frame}{完全シーケンス(2/5)：Authentication}
  \centering
  \begin{tikzpicture}[scale=0.52, every node/.style={font=\tiny}]
    % Actors
    \node[font=\small\bfseries] at (0,11) {UE};
    \node[font=\small\bfseries] at (4,11) {eNB};
    \node[font=\small\bfseries] at (8,11) {S1N2};
    \node[font=\small\bfseries] at (12,11) {AMF/AUSF};

    % Lifelines
    \draw[thick] (0,10.5) -- (0,0);
    \draw[thick] (4,10.5) -- (4,0);
    \draw[thick] (8,10.5) -- (8,0);
    \draw[thick] (12,10.5) -- (12,0);

    % === Authentication Request (DL方向) ===
    \node[font=\scriptsize\bfseries, GoogleGreen] at (10,10.2) {--- Auth Request ---};

    % AMF generates Auth vectors
    \node[draw, fill=GoogleGreen!10, text width=2.3cm, align=center, font=\tiny, anchor=north west] at (12.3,10.0) {
      \textbf{AMF/AUSF}\\
      5G-AKA vectors\\
      RAND, AUTN
    };

    % Auth Request - AMF → S1N2 (NGAP)
    \draw[->, thick, GoogleGreen] (11.9,8.6) -- (8.1,8.6);
    \node[above, font=\tiny] at (10,8.65) {DL NASTransport};

    % S1N2 変換処理
    \node[draw, thick, fill=GoogleYellow!30, text width=2.5cm, align=center, font=\tiny, anchor=north] at (8,8.0) {
      \textbf{NAS変換}\\
      5G Auth → EPS Auth\\
      ngKSI → eKSI
    };

    % Auth Request - S1N2 → eNB (S1AP)
    \draw[->, thick, GoogleBlue] (7.9,6.4) -- (4.1,6.4);
    \node[above, font=\tiny] at (6,6.45) {DL NASTransport};

    % Auth Request - eNB → UE (RRC)
    \draw[->, thick] (3.9,5.6) -- (0.1,5.6);
    \node[above, font=\tiny] at (2,5.65) {AuthenticationRequest};

    % UE側詳細
    \node[draw, fill=GoogleBlue!10, text width=2cm, align=left, font=\tiny, anchor=north east] at (-0.3,5.2) {
      \textbf{4G NAS}\\
      RAND, AUTN\\
      eKSI
    };

    % === Authentication Response (UL方向) ===
    \node[font=\scriptsize\bfseries, GoogleBlue] at (2,4.0) {--- Auth Response ---};

    % UE computes RES
    \node[draw, fill=GoogleBlue!10, text width=2cm, align=center, font=\tiny, anchor=north east] at (-0.3,3.7) {
      \textbf{UE計算}\\
      AUTN検証, RES
    };

    % Auth Response - UE → eNB
    \draw[->, thick] (0.1,2.8) -- (3.9,2.8);
    \node[above, font=\tiny] at (2,2.85) {AuthenticationResponse};

    % Auth Response - eNB → S1N2
    \draw[->, thick, GoogleBlue] (4.1,2.2) -- (7.9,2.2);
    \node[above, font=\tiny] at (6,2.25) {UL NASTransport};

    % RES変換
    \node[draw, thick, fill=GoogleRed!20, text width=2.5cm, align=center, font=\tiny, anchor=north] at (8,1.6) {
      \textbf{RES変換}\\
      RES → RES* (SHA-256)
    };

    % Auth Response - S1N2 → AMF
    \draw[->, thick, GoogleGreen] (8.1,0.4) -- (11.9,0.4);
    \node[above, font=\tiny] at (10,0.45) {UL NASTransport (RES*)};
  \end{tikzpicture}
\end{frame}

%% ========================================
%% シーケンス 3/5: Security Mode（詳細）
%% ========================================
\begin{frame}{完全シーケンス(3/5)：Security Mode}
  \centering
  \begin{tikzpicture}[scale=0.52, every node/.style={font=\tiny}]
    % Actors
    \node[font=\small\bfseries] at (0,11) {UE};
    \node[font=\small\bfseries] at (4,11) {eNB};
    \node[font=\small\bfseries] at (8,11) {S1N2};
    \node[font=\small\bfseries] at (12,11) {AMF};

    % Lifelines
    \draw[thick] (0,10.5) -- (0,0);
    \draw[thick] (4,10.5) -- (4,0);
    \draw[thick] (8,10.5) -- (8,0);
    \draw[thick] (12,10.5) -- (12,0);

    % === Security Mode Command (DL方向) ===
    \node[font=\scriptsize\bfseries, GoogleGreen] at (10,10.2) {--- Security Mode Command ---};

    % AMF generates SMC
    \node[draw, fill=GoogleGreen!10, text width=2.3cm, align=center, font=\tiny, anchor=north west] at (12.3,10.0) {
      \textbf{AMF}\\
      K\_AMF導出\\
      NAS keys生成
    };

    % SMC - AMF → S1N2 (NGAP)
    \draw[->, thick, GoogleGreen] (11.9,8.6) -- (8.1,8.6);
    \node[above, font=\tiny] at (10,8.65) {DL NASTransport (SMC)};

    % SMC詳細 (AMF側)
    \node[draw, fill=GoogleGreen!10, text width=2cm, align=left, font=\tiny, anchor=north west] at (12.3,8.2) {
      \textbf{5G SMC}\\
      Selected algo\\
      ngKSI
    };

    % S1N2 変換処理
    \node[draw, thick, fill=GoogleYellow!30, text width=2.5cm, align=center, font=\tiny, anchor=north] at (8,7.8) {
      \textbf{SMC変換}\\
      5G SMC → EPS SMC\\
      ngKSI → eKSI
    };

    % SMC - S1N2 → eNB (S1AP)
    \draw[->, thick, GoogleBlue] (7.9,6.2) -- (4.1,6.2);
    \node[above, font=\tiny] at (6,6.25) {DL NASTransport};

    % SMC - eNB → UE (RRC)
    \draw[->, thick] (3.9,5.4) -- (0.1,5.4);
    \node[above, font=\tiny] at (2,5.45) {SecurityModeCommand};

    % UE側詳細
    \node[draw, fill=GoogleBlue!10, text width=2cm, align=left, font=\tiny, anchor=north east] at (-0.3,5.0) {
      \textbf{4G SMC}\\
      EIA/EEA\\
      eKSI
    };

    % === Security Mode Complete (UL方向) ===
    \node[font=\scriptsize\bfseries, GoogleBlue] at (2,3.8) {--- Security Mode Complete ---};

    % UE processes SMC
    \node[draw, fill=GoogleBlue!10, text width=2cm, align=center, font=\tiny, anchor=north east] at (-0.3,3.5) {
      \textbf{UE処理}\\
      K\_eNB導出\\
      NAS暗号化開始
    };

    % SMC Complete - UE → eNB
    \draw[->, thick] (0.1,2.4) -- (3.9,2.4);
    \node[above, font=\tiny] at (2,2.45) {SecurityModeComplete};

    % SMC Complete - eNB → S1N2
    \draw[->, thick, GoogleBlue] (4.1,1.8) -- (7.9,1.8);
    \node[above, font=\tiny] at (6,1.85) {UL NASTransport};

    % SMC Complete - S1N2 → AMF
    \draw[->, thick, GoogleGreen] (8.1,1.2) -- (11.9,1.2);
    \node[above, font=\tiny] at (10,1.25) {UL NASTransport (SMC Cmp)};

    % 完了状態
    \node[draw, fill=GoogleGreen!30, text width=3cm, align=center, font=\tiny, anchor=north west] at (12.3,0.8) {
      NAS Security確立完了
    };
  \end{tikzpicture}
\end{frame}

\begin{frame}{完全シーケンス(4/5)：Registration Accept + Early S1 Setup}
  \centering
  \begin{tikzpicture}[scale=0.52, every node/.style={font=\tiny}]
    % Actors
    \node[font=\small\bfseries] at (0,10) {UE};
    \node[font=\small\bfseries] at (4,10) {eNB};
    \node[font=\small\bfseries] at (8,10) {S1N2};
    \node[font=\small\bfseries] at (12,10) {AMF};

    % Lifelines
    \draw[thick] (0,9.5) -- (0,0);
    \draw[thick] (4,9.5) -- (4,0);
    \draw[thick] (8,9.5) -- (8,0);
    \draw[thick] (12,9.5) -- (12,0);

    % === Registration Accept ===
    \node[font=\scriptsize\bfseries, GoogleGreen] at (10,9.2) {--- Registration Accept ---};

    % Reg Accept - AMF → S1N2
    \draw[->, thick, GoogleGreen] (11.9,8.7) -- (8.1,8.7);
    \node[above, font=\tiny] at (10,8.8) {DL NASTransport (Reg Accept)};

    % === Early S1 Setup (独自シーケンス) ===
    \node[font=\scriptsize\bfseries, GoogleRed] at (6,7.8) {--- Early S1 Setup ---};

    % S1N2がインターセプト (右端に配置)
    \node[draw, thick, fill=GoogleRed!20, text width=2.3cm, align=center, font=\tiny, anchor=north west] at (12.3,8.0) {
      PDU Session前に\\
      S1ベアラ先行確立
    };

    % InitialContextSetupRequest (S1) - S1N2 → eNB
    \draw[->, thick, GoogleBlue] (7.9,6.8) -- (4.1,6.8);
    \node[above, font=\tiny] at (6,6.9) {InitialContextSetupReq};

    % E-RAB詳細 (eNBの左側、矢印より下)
    \node[draw, fill=GoogleBlue!10, text width=2cm, align=left, font=\tiny, anchor=north east] at (3.7,6.5) {
      \textbf{E-RAB Setup}\\
      S1N2 Proxy TEID
    };

    % RRC Reconfig - eNB → UE
    \draw[->, thick] (3.9,4.8) -- (0.1,4.8);
    \node[above, font=\tiny] at (2,4.9) {RRC Reconfig + Attach Acc};

    % 4G NAS詳細 (左端、矢印より下)
    \node[draw, fill=GoogleBlue!10, text width=1.8cm, align=left, font=\tiny, anchor=north east] at (-0.3,4.5) {
      \textbf{Attach Acc}\\
      GUTI\\
      Act Def Bearer
    };

    % RRC Reconfig Cmp - UE → eNB
    \draw[->, thick] (0.1,2.8) -- (3.9,2.8);
    \node[above, font=\tiny] at (2,2.9) {RRC Reconfig Cmp};

    % InitialContextSetupResponse (S1) - eNB → S1N2
    \draw[->, thick, GoogleBlue] (4.1,2.0) -- (7.9,2.0);
    \node[above, font=\tiny] at (6,2.1) {InitialContextSetupRes};

    % eNB TEID詳細 (中央、矢印より下)
    \node[draw, fill=GoogleYellow!30, text width=2cm, align=center, font=\tiny, anchor=north] at (8,1.6) {
      eNB TEID 保存
    };
  \end{tikzpicture}
\end{frame}

\begin{frame}{完全シーケンス(5/5)：PDU Session + TEID更新}
  \centering
  \begin{tikzpicture}[scale=0.52, every node/.style={font=\tiny}]
    % Actors
    \node[font=\small\bfseries] at (0,10) {UE};
    \node[font=\small\bfseries] at (4,10) {eNB};
    \node[font=\small\bfseries] at (8,10) {S1N2};
    \node[font=\small\bfseries] at (12,10) {AMF/UPF};

    % Lifelines
    \draw[thick] (0,9.5) -- (0,0.5);
    \draw[thick] (4,9.5) -- (4,0.5);
    \draw[thick] (8,9.5) -- (8,0.5);
    \draw[thick] (12,9.5) -- (12,0.5);

    % === Attach/Registration Complete ===
    \node[font=\scriptsize\bfseries, GoogleBlue] at (6,9.2) {--- Complete + PDU Session ---};

    % Attach Complete - UE → eNB
    \draw[->, thick] (0.1,8.7) -- (3.9,8.7);
    \node[above, font=\tiny] at (2,8.8) {Attach Complete};

    % UL NAS - eNB → S1N2
    \draw[->, thick, GoogleBlue] (4.1,8.1) -- (7.9,8.1);
    \node[above, font=\tiny] at (6,8.2) {UL NASTransport};

    % NAS変換 (右端に配置)
    \node[draw, fill=GoogleYellow!30, text width=2.3cm, align=center, font=\tiny, anchor=north west] at (8.3,7.8) {
      → Reg Complete\\
      + PDU Session Req
    };

    % Reg Complete - S1N2 → AMF
    \draw[->, thick, GoogleGreen] (8.1,6.4) -- (11.9,6.4);
    \node[above, font=\tiny] at (10,6.5) {Registration Complete};

    % PDU Session Req - S1N2 → AMF
    \draw[->, thick, GoogleGreen] (8.1,5.6) -- (11.9,5.6);
    \node[above, font=\tiny] at (10,5.7) {PDU Session Est Req};

    % === PDU Session Resource Setup ===
    \node[font=\scriptsize\bfseries, GoogleGreen] at (10,4.8) {--- PDU Session Setup ---};

    % PDU Res Setup - AMF → S1N2
    \draw[->, thick, GoogleGreen] (11.9,4.2) -- (8.1,4.2);
    \node[above, font=\tiny] at (10,4.3) {PDUSessionResSetupReq};

    % UPF TEID詳細 (右端、矢印より下)
    \node[draw, fill=GoogleGreen!10, text width=2cm, align=left, font=\tiny, anchor=north west] at (12.3,3.9) {
      \textbf{PDU Session}\\
      UPF N3 TEID
    };

    % TEID更新 (中央、矢印より下)
    \node[draw, thick, fill=GoogleRed!20, text width=2.3cm, align=center, font=\tiny, anchor=north] at (8,3.6) {
      \textbf{TEID更新}\\
      Proxy → UPF TEID
    };

    % PDU Res Setup Res - S1N2 → AMF
    \draw[->, thick, GoogleGreen] (8.1,2.0) -- (11.9,2.0);
    \node[above, font=\tiny] at (10,2.1) {PDUSessionResSetupRes};

    % 最終状態
    \node[draw, fill=GoogleGreen!30, text width=4cm, align=center, font=\tiny] at (6,1.0) {
      \textbf{接続完了} - GTP-U TEID変換によりデータ疎通
    };
  \end{tikzpicture}
\end{frame}

\begin{frame}{NASメッセージ変換対応表}
  \begin{columns}[T]
    \begin{column}{0.48\textwidth}
      \textbf{上り (UE → Core)}
      \begin{table}
        \tiny
        \begin{tabular}{|l|l|}
          \hline
          \textbf{4G NAS} & \textbf{5G NAS} \\
          \hline
          Attach Request & Registration Request \\
          \hline
          Auth Response & Auth Response \\
          (RES) & (RES*: HRES*計算) \\
          \hline
          Security Mode Cmp & Security Mode Cmp \\
          \hline
          Attach Complete & Registration Complete \\
          + ESM情報 & + PDU Session Est Req \\
          \hline
        \end{tabular}
      \end{table}

      \vspace{0.3cm}
      \textbf{ID変換}
      \begin{table}
        \tiny
        \begin{tabular}{|l|l|}
          \hline
          \textbf{4G} & \textbf{5G} \\
          \hline
          IMSI & SUCI (NULL scheme) \\
          \hline
          GUTI & 5G-GUTI \\
          \hline
          eKSI (3bit) & ngKSI (4bit, TSC=0) \\
          \hline
        \end{tabular}
      \end{table}
    \end{column}
    \begin{column}{0.48\textwidth}
      \textbf{下り (Core → UE)}
      \begin{table}
        \tiny
        \begin{tabular}{|l|l|}
          \hline
          \textbf{5G NAS} & \textbf{4G NAS} \\
          \hline
          Auth Request & Auth Request \\
          (ngKSI) & (eKSI) \\
          \hline
          Security Mode Cmd & Security Mode Cmd \\
          (NEA/NIA) & (EEA/EIA) \\
          \hline
          Registration Accept & Attach Accept \\
          + 5G-GUTI & + GUTI \\
          \hline
          PDU Session Est Acc & Act Def EPS Bearer \\
          (5QI, DNN) & (QCI, APN) \\
          \hline
        \end{tabular}
      \end{table}

      \vspace{0.3cm}
      \textbf{セキュリティアルゴリズム}
      \begin{table}
        \tiny
        \begin{tabular}{|l|l|}
          \hline
          \textbf{4G} & \textbf{5G} \\
          \hline
          EEA0/EIA0 & NEA0/NIA0 \\
          EEA1/EIA1 & NEA1/NIA1 \\
          EEA2/EIA2 & NEA2/NIA2 \\
          \hline
        \end{tabular}
      \end{table}
    \end{column}
  \end{columns}
\end{frame}

\begin{frame}{データプレーン：GTP-U TEID変換}
  \centering
  \begin{tikzpicture}[scale=0.8, every node/.style={font=\small}]
    % Actors
    \node[font=\normalsize\bfseries] at (0,5) {eNB};
    \node[font=\normalsize\bfseries] at (6,5) {S1N2 Converter};
    \node[font=\normalsize\bfseries] at (12,5) {UPF};

    % Lifelines
    \draw[thick] (0,4.5) -- (0,0);
    \draw[thick] (6,4.5) -- (6,0);
    \draw[thick] (12,4.5) -- (12,0);

    % Uplink
    \draw[->, thick, GoogleBlue] (0,3.8) -- (6,3.6);
    \node[above, font=\scriptsize, align=center] at (3,3.8) {GTP-U\\S1-U TEID};

    \draw[->, thick, GoogleGreen] (6,3.3) -- (12,3.1);
    \node[above, font=\scriptsize, align=center] at (9,3.3) {GTP-U\\N3 TEID};

    % Downlink
    \draw[<-, thick, GoogleGreen] (6,2.3) -- (12,2.1);
    \node[below, font=\scriptsize, align=center] at (9,2.1) {GTP-U\\N3 TEID};

    \draw[<-, thick, GoogleBlue] (0,1.8) -- (6,1.6);
    \node[below, font=\scriptsize, align=center] at (3,1.6) {GTP-U\\S1-U TEID};

    % TEID Mapping Table
    \node[draw, thick, fill=GoogleYellow!20, text width=4cm, align=left, font=\scriptsize] at (6,0.3) {
      \textbf{TEID Mapping Table}\\
      S1-U: 0x12345678\\
      $\updownarrow$\\
      N3: 0xABCDEF00
    };
  \end{tikzpicture}
\end{frame}

\begin{frame}{主要な変換パラメータ}
  \begin{columns}[T]
    \begin{column}{0.48\textwidth}
      \textbf{C-Plane変換}
      \begin{itemize}
        \item \textbf{ID変換}
          \begin{itemize}
            \item ENB-ID $\rightarrow$ gNB-ID
            \item IMSI $\rightarrow$ SUCI
            \item GUTI $\rightarrow$ 5G-GUTI
          \end{itemize}
        \item \textbf{メッセージ変換}
          \begin{itemize}
            \item S1SetupReq $\rightarrow$ NGSetupReq
            \item AttachReq $\rightarrow$ RegistrationReq
          \end{itemize}
      \end{itemize}
    \end{column}
    \begin{column}{0.48\textwidth}
      \textbf{U-Plane変換}
      \begin{itemize}
        \item \textbf{QoS変換}
          \begin{itemize}
            \item QCI $\rightarrow$ 5QI
            \item E-RAB ID $\rightarrow$ PDU Session ID
          \end{itemize}
        \item \textbf{TEID管理}
          \begin{itemize}
            \item S1-U TEID $\leftrightarrow$ N3 TEID
            \item マッピングテーブル管理
          \end{itemize}
      \end{itemize}
    \end{column}
  \end{columns}

  \vspace{0.5cm}
  \begin{block}{実装のポイント}
    \begin{itemize}
      \item UE毎のコンテキスト管理（ENB-UE-ID、AMF-UE-ID、TEID）
      \item NASメッセージの透過的な転送（暗号化されたまま変換）
      \item セキュリティパラメータの保持（eKSI $\leftrightarrow$ ngKSI）
    \end{itemize}
  \end{block}
\end{frame}

\section{実験・評価}
\begin{frame}{実験環境：sXGPとOpen5GSを用いた実機検証}
  \centering
  \begin{tikzpicture}[scale=0.85, every node/.style={font=\small}]
    % UE
    \node[draw, rounded corners, fill=GoogleLightGray, minimum width=1.8cm, minimum height=0.8cm, align=center] (ue) at (0,0) {UE\\{\scriptsize Pixel}};

    % sXGP
    \node[draw, rounded corners, fill=GoogleBlue!30, minimum width=2cm, minimum height=0.8cm, align=center] (sxgp) at (3.5,0) {sXGP AP\\{\scriptsize 1.9GHz}};

    % Gateway
    \node[draw, thick, rounded corners, fill=GoogleYellow!20, minimum width=2.5cm, minimum height=1cm, align=center] (gw) at (7.5,0) {Gateway\\{\scriptsize C++/Python}};

    % Open5GS
    \node[draw, rounded corners, fill=GoogleGreen!30, minimum width=2cm, minimum height=0.8cm, align=center] (o5gs) at (11.5,0) {Open5GS\\{\scriptsize 5GC}};

    % Connections
    \draw[->, thick] (ue) -- (sxgp) node[midway, above] {\scriptsize LTE};
    \draw[->, thick] (sxgp) -- (gw) node[midway, above] {\scriptsize S1};
    \draw[->, thick] (gw) -- (o5gs) node[midway, above] {\scriptsize N2/N3};

    % Annotations
    \node[below=0.3cm of ue, text width=1.8cm, align=center, font=\tiny] {4G/5G\\Dual SIM};
    \node[below=0.3cm of sxgp, text width=2cm, align=center, font=\tiny] {プライベートLTE\\(4G eNB)};
    \node[below=0.5cm of gw, text width=2.5cm, align=center, font=\tiny] {提案実装\\疎結合変換層};
    \node[below=0.3cm of o5gs, text width=2cm, align=center, font=\tiny] {オープンソース\\5GC実装};
  \end{tikzpicture}

  \vspace{0.5cm}
  \begin{itemize}
    \item \textbf{検証目的:} 4G無線経由で5GC登録とデータ通信が成立するか
    \item \textbf{評価軸:} ①機能検証(定性)、②性能評価(定量)
  \end{itemize}
\end{frame}

\begin{frame}{評価結果(1)：C-Plane接続検証}
  \begin{itemize}
    \item \textbf{結果：Registration成功}
      \begin{itemize}
        \item 4G経由で送られた \texttt{Attach Request} 内の5G NASメッセージを抽出・転送
        \item 5GC側で正常に \texttt{Registration Accept} が発行され、PDU Sessionが確立
      \end{itemize}
    \item \textbf{検証の意義 (定性評価)}
      \begin{itemize}
        \item \textbf{依存の排除:} ゲートウェイが差異を吸収し、直接接続不可なRAN-Core間の接続を実現
        \item \textbf{機能の抽象化:} 世代共通のAAA/RM(Registration Management)機能を維持したままマッピングに成功
      \end{itemize}
  \end{itemize}
\end{frame}

\begin{frame}{評価結果(2)：U-Plane性能評価}
  \begin{itemize}
    \item \textbf{測定の狙い}
      \begin{itemize}
        \item \textbf{相互接続の実証:} Ping/iPerfの疎通成功により、プロトコル変換の正常性を証明
        \item \textbf{実用性の評価:} ゲートウェイ有無(既存LTE vs 提案方式)の比較によりオーバーヘッドを検証
      \end{itemize}
    \item \textbf{結果概要}
      \begin{itemize}
        \item \textbf{遅延:} 処理による追加遅延は数ms程度
        \item \textbf{スループット:} 4G無線区間の帯域がボトルネックとなり、変換処理による低下は軽微
      \end{itemize}
    \item \textbf{考察}
      \begin{itemize}
        \item 変換処理のオーバーヘッドは無線区間の遅延・帯域に対して十分小さい
        \item $\Rightarrow$ \textbf{異世代相互接続は性能面でも実用可能}であることを確認
      \end{itemize}
  \end{itemize}
\end{frame}

\section{結論}
\begin{frame}{学術的貢献と実用的インパクト}
  \begin{block}{学術的貢献}
    \begin{itemize}
      \item 世代不変性に基づく\alert{疎結合アーキテクチャ}の理論化と実証
      \item RANとCoreの進化サイクルを分離する新しい移行モデルの提示
    \end{itemize}
  \end{block}
  \begin{block}{実用的インパクト}
    \begin{itemize}
      \item \textbf{R\&D:} 低コストな次世代Core検証環境(既存設備の再利用)
      \item \textbf{商用展開:} RAN更新を伴わない柔軟なCore先行移行戦略
      \item \textbf{NTN/IoT:} 物理更新不可能なインフラの延命と機能拡張
    \end{itemize}
  \end{block}
\end{frame}

\begin{frame}{今後の展望}
  \begin{itemize}
    \item \textbf{スライシング制御の実装}
      \begin{itemize}
        \item レガシーRAN配下の端末に対するCore側QoS制御の実証
      \end{itemize}
    \item \textbf{NTN環境での遅延耐性評価}
      \begin{itemize}
        \item 衛星/HAPS特有の長遅延(数百ms)下でのプロトコル動作検証
      \end{itemize}
    \item \textbf{6G Coreへの拡張}
      \begin{itemize}
        \item 提案手法の汎用性検証(4G→5Gだけでなく5G→6Gへも適用可能か)
      \end{itemize}
  \end{itemize}
\end{frame}

\end{document}
